\section{Exact binary sumsets}

Which integers are sums of two primitive values of a quadratic norm?
For $x^2+2y^2$, divisibility by sixteen is the complete obstruction.
For $x^2+xy+2y^2$, the local obstruction is divisibility by forty-nine,
but three additional integers fail.  This paper determines the full
sumsets in these examples and in two further class-number-one cases.
The distinction between eventual existence and a complete classification
is the point: every small exception must be proved, and every remaining
finite range must meet an infinite argument.

For a negative discriminant $d$, put
\[
 P_d=\{Q_d(x,y)>0:x,y\in\Z,\ \gcd(x,y)=1\},\qquad
 \begin{array}{c|c}
 d&Q_d(x,y)\\\hline
 -3&x^2-xy+y^2\\
 -7&x^2+xy+2y^2\\
 -8&x^2+2y^2\\
 -11&x^2+xy+3y^2.
 \end{array}
\]
Zero coordinates and both signs are allowed.  Each summand is a positive
numerical value; different coordinate representations do not create
new summands.

\begin{theorem}\label{thm:exact-overview}
The exact binary sumsets are
\[
\begin{array}{c|l}
 d&P_d+P_d\\\hline
 -3&\{N\ge2:N\text{ even},\ N\not\equiv0,12\pmod{18}\},\\
 -7&\{N\ge2:49\nmid N\}\setminus\{7,77,203\},\\
 -8&\{N\ge2:16\nmid N\},\\
 -11&\{N\ge2:N\text{ even},\ 121\nmid N\}.
\end{array}
\]
\end{theorem}

For the Eisenstein row, let
\[
 \AA=\{n\ge1:p\mid n\Longrightarrow p\equiv1\pmod3\}.
\]
Here $1\in\AA$.  The elementary identity
\[
                   P_{-3}=\AA\sqcup(3\!\cdot\!\AA)
\]
explains why two split-prime core identities suffice.  An inert prime
cannot divide a primitive norm, while the ramified prime three occurs
at most once; conversely every value in either displayed part has
coprime coordinates.  Proposition~\ref{prop:primitive-classification}
proves this in the arithmetic preliminaries.  Thus the three types of
binary sums are $\AA+\AA$, $\AA+3\!\cdot\!\AA$ and
$3\!\cdot\!(\AA+\AA)$.

The proofs establish stronger identities for the semigroups generated
by split primes.  In each ring, a primitive norm has no inert prime
factor, and the ramified prime can appear only once.  Thus the norm
value set is a split-prime semigroup together with one ramified
dilation.  The residue restrictions are elementary; proving their
sufficiency is a binary additive problem.

Our method adapts the quaternionic norm-shell and torus-normalizer
argument used for the Gaussian restricted-prime theorem of Howe,
Huffman, Kohlmeyer and Kruer \cite[Theorem~1.1]{Euler}.  A quaternion
shell produces many candidate pairs of binary norms.  The Hecke
spectrum bounds the fraction spoiled by each forbidden prime.
Exact dyadic and ramified filters keep the total fraction below one.
The spectral input is classical Jacquet--Langlands and the weight-two
Ramanujan bound; Section~\ref{sec:exact-hecke} states its normalization.
The applications require different local quotients, including all ideal
classes and inverse-unit mass at discriminant eleven.

Explicit estimates settle an infinite range.  Rational interval
certificates bridge that range down to complete finite sumset scans;
multiplicative descent then closes all remaining targets.  A finite
scan by itself does not imply any of the unrestricted identities.
The supplied programs reconstruct the membership arrays, verify the
interval inequalities in exact arithmetic, and test every target in
the stated finite domains.  Their mathematical roles and completeness
arguments are included below and summarized in
Appendix~\ref{app:verifier}.

Table~\ref{tab:exact-coverage} shows how the finite and infinite parts
meet.  In its last three rows, $C_d$ is the split-prime core defined in
the corresponding section.  The variable $n$ is a reduced target; in
the $-8$ row, $m$ is an odd inert-supported core.  Interval certificates
certify the stated families throughout each interval, rather than
scanning its individual targets.

\begin{table}[ht]
\centering\small
\setlength{\tabcolsep}{3pt}
\renewcommand{\arraystretch}{1.13}
\caption{Coverage of the exact classifications.  All finite scans include
only the locally eligible targets of the indicated families.}
\label{tab:exact-coverage}
\begin{tabular}{@{}>{\raggedright\arraybackslash}p{.045\linewidth}>{\raggedright\arraybackslash}p{.17\linewidth}>{\raggedright\arraybackslash}p{.255\linewidth}>{\raggedright\arraybackslash}p{.225\linewidth}>{\raggedright\arraybackslash}p{.24\linewidth}@{}}
\toprule
$d$&Complete finite scan&Rational interval bridge&Infinite estimate&Descent and dilation\\
\midrule
$-3$&$N<10^8$: $\AA+\AA$ and $\AA+3\!\cdot\!\AA$.
&No interval bridge needed.
&Reduced $n\ge10^8$, in both admissible classes modulo six.
&Remove all split-prime factors, then scale back; use $P_{-3}=\AA\sqcup3\!\cdot\!\AA$.\\[3pt]
$-7$&$N\le2^{24}$: $C_7+C_7$ and $P_7+P_7$, with their listed exceptions.
&$257$ intervals on $[2^{24},2^{37})$; $17$ on $[2^{37},2^{54})$, for inert-supported targets and split primes.
&The same reduced families for $n\ge2^{54}$.
&Remove split factors; patch quotients $11,29$.  Ramified dilation by seven gives the primitive exceptions.\\[3pt]
$-8$&$N\le2^{27}$: $C_8+C_8$ and $C_8+2C_8$.
&$107$ intervals: $m\in[2^{27},2^{60})$ for mixed sums, $[2^{26},2^{60})$ at $2m$, and $[2^{20},2^{60})$ at $4m$.
&Inert cores $m\ge2^{60}$; the mixed prime boundary holds for primes $N\ge2^{22}$.
&Remove split factors; the remaining targets are $m,2m,4m$.  Use $P_{-8}=C_8\sqcup2C_8$.\\[3pt]
$-11$&$N\le2^{25}$: $C_{11}+C_{11}$.
&$35$ intervals on $[2^{25},2^{60})$, for even inert-supported targets prime to eleven.
&The same reduced family for $n\ge2^{60}$.
&Remove split factors; ramified dilation by eleven gives $P_{-11}+P_{-11}$.\\
\bottomrule
\end{tabular}
\end{table}

We first prove the Eisenstein identities, including the mixed family
forced by the ramified prime three.  The following sections treat
$-7,-8,-11$, and explain how their finite certificates meet their
respective infinite estimates.  The exact exceptions at $-7$ are the
only exceptions in Theorem~\ref{thm:exact-overview} beyond its displayed
congruence conditions.
